\documentclass[11pt,a4paper]{article}
\usepackage[utf8]{inputenc}
\usepackage{geometry}
\geometry{margin=1in}
\usepackage{lineno}
\usepackage{xcolor}
\usepackage{titlesec}
\usepackage{array}

% Command for redacting personal information with a black box
\newcommand{\redacted}{\rule{2.5cm}{2ex}}

\title{\textbf{Social Science Fair Interview 02 \\ Transcript: Value and Pressure of Football Defenders (Winger Perspective)}}
\author{\textbf{Research Team Members:} Luke Sun, Kiara Wei, Winston Huang}
\date{May 28, 2026}

\begin{document}

\maketitle

\section*{Transcript Metadata}
\begin{tabular}{ll}
    \textbf{Project Theme:} & Discussion on the Value and Pressure of Football Defenders \\
    \textbf{Date of Interview:} & Thursday, May 28, 2026 \\
    \textbf{Time of Interview:} & 20:09 -- 20:16 (GMT+08) \\
    \textbf{Interviewer Identifier:} & Speaker 2 Luke Sun\\
    \textbf{Interviewee Identifier:} & Speaker 1 (\redacted) \\
    \textbf{Interviewee Subject Profile:} & School Team Athlete (Left-Winger)
\end{tabular}

\vspace{2em}
\hrule
\vspace{1.5em}

\section*{Official Verbatim Transcript}

\begin{linenumbers}

\noindent \textbf{Speaker 1 [00:00]:} You are about to cross it.

\vspace{0.5em}
\noindent \textbf{Speaker 2 [00:01]:} You are \redacted, correct? 

\vspace{0.5em}
\noindent \textbf{Speaker 1 [00:02]:} Yes, I am.

\vspace{0.5em}
\noindent \textbf{Speaker 2 [00:03]:} Okay. What position do you currently play on the school football team?

\vspace{0.5em}
\noindent \textbf{Speaker 1 [00:10]:} Left-winger.

\vspace{0.5em}
\noindent \textbf{Speaker 2 [00:13]:} Great. How do you define your core responsibility on the pitch? What do you primarily do?

\vspace{0.5em}
\noindent \textbf{Speaker 1 [00:20]:} I am responsible for wing penetration, ball circulation, as well as assisting my teammates or scoring myself.

\vspace{0.5em}
\noindent \textbf{Speaker 2 [00:30]:} Mhm, okay. In the last match, how much of a decisive role do you think you played in the team's victory or defeat?

\vspace{0.5em}
\noindent \textbf{Speaker 1 [00:37]:} When was the last match?

\vspace{0.5em}
\noindent \textbf{Speaker 2 [00:40]:} The league match, the previous one.

\vspace{0.5em}
\noindent \textbf{Speaker 1 [00:46]:} Unfortunately, not much; I only managed to get one assist, which is a pity. However, it was a highly critical assist, so I believe it still played a key role.

\vspace{0.5em}
\noindent \textbf{Speaker 2 [01:03]:} Key in which aspect? Defensive or offensive?

\vspace{0.5em}
\noindent \textbf{Speaker 1 [01:06]:} Creating a threat during the offense.

\vspace{0.5em}
\noindent \textbf{Speaker 2 [01:09]:} Pressing? 

\vspace{0.5em}
\noindent \textbf{Speaker 1 [01:10]:} Yes.

\vspace{0.5em}
\noindent \textbf{Speaker 2 [01:11]:} Okay, during the match, if in the early stages---say, the first few minutes---you miss a golden opportunity, but later in the match, you score a goal. Do you think people will still hold your previous mistake against you?

\vspace{0.5em}
\noindent \textbf{Speaker 1 [01:27]:} No, they won't.

\vspace{0.5em}
\noindent \textbf{Speaker 2 [01:30]:} When you see backfield players---such as defenders, or people like \redacted---make a mistake, like misplacing a pass or missing a tackle, leading directly to conceding a goal, what kind of emotional reaction do you typically have?

\vspace{0.5em}
\noindent \textbf{Speaker 1 [01:47]:} It doesn't matter; as long as we didn't concede too many goals, it's fine.

\vspace{0.5em}
\noindent \textbf{Speaker 2 [01:54]:} Even if it's a direct error, like a severe miss or letting the ball slip through, it still doesn't matter?

\vspace{0.5em}
\noindent \textbf{Speaker 1 [02:02]:} My mindset might be somewhat affected.

\vspace{0.5em}
\noindent \textbf{Speaker 2 [02:06]:} What specific kind of impact would that be?

\vspace{0.5em}
\noindent \textbf{Speaker 1 [02:09]:} Possibly, if it leads to being equalized or overturned, I might feel some psychological pressure in the subsequent offense, which could prevent me from executing passing and control more effectively.

\vspace{0.5em}
\noindent \textbf{Speaker 2 [02:26]:} Understood. Do you agree with the statement that the cost of making a mistake for a defensive player is far higher than that for an attacking player?

\vspace{0.5em}
\noindent \textbf{Speaker 1 [02:35]:} Yes, I believe so.

\vspace{0.5em}
\noindent \textbf{Speaker 2 [02:36]:} Could you elaborate on the reasons for that?

\vspace{0.5em}
\noindent \textbf{Speaker 1 [02:38]:} For instance, when an attacking player loses the ball, they don't bear a significant amount of pressure because they are not the first line---there are still people behind to back them up. However, if a defender loses the ball, they are the final line of defense; if they lose it, the opponent's probability of scoring directly becomes much higher.

\vspace{0.5em}
\noindent \textbf{Speaker 2 [02:59]:} Mhm. Do you feel that goals and assists grant you greater exposure on campus?

\vspace{0.5em}
\noindent \textbf{Speaker 1 [03:05]:} Yes, they do.

\vspace{0.5em}
\noindent \textbf{Speaker 2 [03:07]:} Compared to metrics like a defender's critical clearances, do goals and assists carry greater topicality or generate more discussion than clearances or tackles?

\vspace{0.5em}
\noindent \textbf{Speaker 1 [03:24]:} I think so. Most highlight reels focus on brilliant offensive moments; there are rarely scenes showcasing outstanding defensive plays.

\vspace{0.5em}
\noindent \textbf{Speaker 2 [03:38]:} Do you believe this extra attention is purely based on the aesthetic appeal or skill of the goal itself, or because scoring inherently possesses greater topicality than tackles and clearances?

\vspace{0.5em}
\noindent \textbf{Speaker 1 [03:52]:} I think it leans a bit more toward aesthetic appeal.

\vspace{0.5em}
\noindent \textbf{Speaker 2 [03:55]:} Meaning the aesthetic value of a goal is slightly higher than that of a sliding tackle or defense? Right, right. Then do you think this somewhat unbalanced return could have adverse effects on backfield players? Such as inducing a psychological state where backfield players feel neglected?

\vspace{0.5em}
\noindent \textbf{Speaker 1 [04:25]:} Yes, I believe it does.

\vspace{0.5em}
\noindent \textbf{Speaker 2 [04:27]:} Since you just mentioned the asymmetry of risk between both sides---where frontfield players face lower error costs and backfield players face higher ones---do you think there is a way to resolve this current situation? For instance, editing more defensive highlights online to promote the art of defending. Or, during daily training or football classrooms, advocating more for a collective spirit, discouraging individual heroism, and instead embracing offensive-defensive balance or even placing greater emphasis on defense. Do you think there are similar solutions or measures? Can you think of any?

\vspace{0.5em}
\noindent \textbf{Speaker 1 [05:10]:} It could indeed make offense and defense more balanced. It is possible that a team is defense-oriented, meaning their progress or performance in defense might be superior to their offense. However, there is also an issue, which is the general public psychology: people tend to perceive offense as more aesthetically pleasing. Their interest in defense is far less than in offense. Therefore, to satisfy media demands, there will probably still be more offensive highlights. Yet within the team, coaches can absolutely place greater emphasis on defense, such as implementing counter-attacking strategies, so that defense would also be highlighted and integrated thoroughly.

\vspace{0.5em}
\noindent \textbf{Speaker 2 [05:54]:} As for the general public, do you think there is any way to slightly alleviate their preference for offense, or is offense inherently characterized by high aesthetic appeal?

\vspace{0.5em}
\noindent \textbf{Speaker 1 [06:13]:} If the public consists of avid fans, or people who possess a deep understanding of the team, they might focus more on the tactics rather than solely on the offensive aspect. Therefore, simply promoting football more broadly to the general public to deepen their understanding could potentially lead them to have higher expectations and appreciation for defense.

\vspace{0.5em}
\noindent \textbf{Speaker 2 [06:43]:} Essentially, elevating their overall comprehension of football as a sport. Yes. Through methods like online promotion, campaigns, public opinion, etc., correct? Yes. Great, that concludes it.

\end{linenumbers}

\end{document}